\documentclass[11pt,a4paper]{article}

\usepackage[utf8]{inputenc}
\usepackage[T1]{fontenc}
\usepackage{geometry}
\usepackage{hyperref}
\usepackage{enumitem}
\usepackage{xcolor}
\usepackage{fancyhdr}
\usepackage{amsmath}
\usepackage{booktabs}
\usepackage{listings}

\geometry{margin=1in}
\hypersetup{
  colorlinks=true,
  linkcolor=blue,
  urlcolor=blue
}

\setlist{itemsep=0.35em, topsep=0.35em}

\pagestyle{fancy}
\fancyhf{}
\lhead{Object Oriented Programming (B.Tech CSE)}
\rhead{Units 01--02 + Unit 03 (Partial): Diagnostic Assessment}
\cfoot{\thepage}
\setlength{\headheight}{14pt}

\lstset{
  language=Java,
  basicstyle=\ttfamily\small,
  keywordstyle=\color{blue},
  commentstyle=\color{gray},
  stringstyle=\color{teal},
  showstringspaces=false,
  columns=fullflexible,
  frame=single
}

\newcommand{\SubmissionDeadline}{March 8, 2026}

% Marks per question (total = 100)
\newcommand{\EMarks}{\textbf{[2 marks]} }
\newcommand{\MMarks}{\textbf{[3 marks]} }
\newcommand{\HMarks}{\textbf{[5 marks]} }

\begin{document}

\begin{center}
  {\LARGE \textbf{Diagnostic Assessment -- Units 01 \& 02 + Unit 03 (Partial)}}\\[0.25em]
  {\large Object Oriented Programming (CSEG1044)}\\[0.25em]
  \normalsize (Unit 01--02 + Unit 03 partial: nested classes + exception handling/handlers; designed to identify learning gaps)
\end{center}

\noindent
\textbf{Instructor:} Tofik Ali \hfill \textbf{Time Limit:} None\\
\textbf{Submission Deadline:} \SubmissionDeadline\\
\textbf{Student Name:} \_\_\_\_\_\_\_\_\_\_\_\_\_\_\_\_\_ \hfill \textbf{Roll No.:} \_\_\_\_\_\_\_\_\\

\vspace{0.6em}
\hrule
\vspace{0.8em}

\textbf{Instructions}
\begin{itemize}[leftmargin=*]
  \item \textbf{This is a diagnostic assessment, not a quiz/test.} The goal is to identify what to revise.
  \item \textbf{Total marks: 100.} Easy: 10 $\times$ 2 = 20, Medium: 10 $\times$ 3 = 30, Hard: 10 $\times$ 5 = 50.
  \item \textbf{No time limit.} Submit on or before \textbf{\SubmissionDeadline}.
  \item \textbf{Important (70\% rule): To have your assignment marks counted, you must score at least 70\% (70/100 or more) in this assessment.}
    If you score below 70\%, you may reattempt and resubmit. Your latest submission will be considered.
  \item \textbf{Marking policy (honest attempt):} Marks are deducted only for (i) not attempting a question and (ii) high similarity with other submissions (copying).
    Correctness is \textbf{not} the main focus; your reasoning and effort matter more.
  \item Answer \textbf{all} questions. For code questions, snippets are acceptable unless a full program is requested.
  \item For MCQs, follow the instruction in the question (some are \textbf{select all that apply}).
\end{itemize}

\vspace{0.6em}

\section*{Easy (10 Questions)}
\begin{enumerate}[label=\textbf{E\arabic*.}, leftmargin=*]
  \item \EMarks \textbf{[MCQ -- Select all that apply]} Which are core OOP principles?
    \begin{enumerate}[label=(\Alph*), leftmargin=*]
      \item Abstraction
      \item Encapsulation
      \item Compilation
      \item Inheritance
      \item Polymorphism
      \item Pagination
    \end{enumerate}

  \item \EMarks \textbf{[MCQ]} If a class is declared \texttt{public class Hello}, what must the file name be?
    \begin{enumerate}[label=(\Alph*), leftmargin=*]
      \item \texttt{Hello.java}
      \item \texttt{hello.java}
      \item \texttt{Main.java}
      \item Any name is fine
    \end{enumerate}

  \item \EMarks \textbf{[MCQ]} Which statement is correct?
    \begin{enumerate}[label=(\Alph*), leftmargin=*]
      \item JVM compiles Java source code to bytecode
      \item JDK is required only to run a \texttt{.class} file
      \item JVM executes bytecode
      \item JRE is a compiler
    \end{enumerate}

  \item \EMarks \textbf{[MCQ]} Which of the following is \textbf{not} a primitive type?
    \begin{enumerate}[label=(\Alph*), leftmargin=*]
      \item \texttt{int}
      \item \texttt{double}
      \item \texttt{boolean}
      \item \texttt{String}
    \end{enumerate}

  \item \EMarks \textbf{[Output]} What is printed?
    \begin{lstlisting}
int x = 3;
while (x > 0) {
    System.out.print(x + " ");
    x--;
}
    \end{lstlisting}

  \item \EMarks \textbf{[Subjective]} Why is \texttt{String} immutable in Java? (2--3 lines)

  \item \EMarks \textbf{[MCQ]} Which keyword is used to inherit a class in Java?
    \begin{enumerate}[label=(\Alph*), leftmargin=*]
      \item \texttt{implements}
      \item \texttt{extends}
      \item \texttt{inherits}
      \item \texttt{super}
    \end{enumerate}

  \item \EMarks \textbf{[MCQ -- Select all that apply]} Which are access modifiers in Java?
    \begin{enumerate}[label=(\Alph*), leftmargin=*]
      \item \texttt{private}
      \item \texttt{protected}
      \item \texttt{public}
      \item \texttt{package}
      \item \texttt{default}
    \end{enumerate}

  \item \EMarks \textbf{[MCQ]} Which of the following is a \textbf{checked} exception?
    \begin{enumerate}[label=(\Alph*), leftmargin=*]
      \item \texttt{IOException}
      \item \texttt{NullPointerException}
      \item \texttt{ArithmeticException}
      \item \texttt{IndexOutOfBoundsException}
    \end{enumerate}

  \item \EMarks \textbf{[MCQ -- Select all that apply]} Which are types of nested classes in Java?
    \begin{enumerate}[label=(\Alph*), leftmargin=*]
      \item Static nested class
      \item Inner (non-static) class
      \item Local class
      \item Anonymous class
      \item Global class
    \end{enumerate}
\end{enumerate}

\section*{Medium (10 Questions)}
\begin{enumerate}[label=\textbf{M\arabic*.}, leftmargin=*]
  \item \MMarks \textbf{[Code]} Write a \texttt{Student} class with a parameterized constructor and use \texttt{this} to assign fields.

  \item \MMarks \textbf{[Code]} Write two overloaded methods \texttt{max(int a, int b)} and \texttt{max(int a, int b, int c)}.

  \item \MMarks \textbf{[Code]} Write a snippet that demonstrates a \texttt{static} variable counting created objects.

  \item \MMarks \textbf{[Code]} Write an interface \texttt{Payable} and one implementing class \texttt{CardPayment}. Call \texttt{pay()} using a \texttt{Payable} reference.

  \item \MMarks \textbf{[Code]} Write a base class \texttt{Vehicle} and derived class \texttt{Car}. Override a method and show dynamic dispatch.

  \item \MMarks \textbf{[Code]} Write correct \texttt{package} and \texttt{import} lines for a class in package \texttt{upes.oop.util}.

  \item \MMarks \textbf{[Code]} Write a \texttt{try/catch/finally} block that handles divide-by-zero and always prints \texttt{Done}.

  \item \MMarks \textbf{[Code]} Write one line showing \texttt{throw} and one line showing \texttt{throws} (any suitable example).

  \item \MMarks \textbf{[Code]} Write code to instantiate:
    (a) a static nested class \texttt{Outer.Nested}, and
    (b) an inner class \texttt{Outer.Inner}.

  \item \MMarks \textbf{[Concept]} What is an exception handler? Write 4--6 lines with one example situation.
\end{enumerate}

\section*{Hard (10 Questions)}
\begin{enumerate}[label=\textbf{H\arabic*.}, leftmargin=*]
  \item \HMarks \textbf{[Program]} Write a complete program that reads two integers and performs division.
    Handle invalid input and divide-by-zero. Use \texttt{finally} to print \texttt{Done}. Include sample output.

  \item \HMarks \textbf{[Program]} Implement \texttt{BankAccount} with \texttt{deposit()} and \texttt{withdraw()}.
    Throw a custom \texttt{InsufficientBalanceException} for invalid withdrawal. Include sample run.

  \item \HMarks \textbf{[Program]} Create an abstract class \texttt{Shape} with method \texttt{area()} and implement \texttt{Circle} and \texttt{Rectangle}.
    Store objects in \texttt{Shape[]} and print all areas (dynamic dispatch).

  \item \HMarks \textbf{[Program]} Create \texttt{InvalidMarksException extends Exception}.
    Write method \texttt{setMarks(int m) throws InvalidMarksException} that validates 0--100.
    Demonstrate using \texttt{try/catch} in \texttt{main}.

  \item \HMarks \textbf{[Program]} Show exception propagation using \texttt{main() -> a() -> b()} where \texttt{b()} throws.
    Catch only in \texttt{main} and print flow messages.

  \item \HMarks \textbf{[Program]} Demonstrate access modifiers across packages:
    create two packages and show what can/cannot be accessed for \texttt{private}, default, \texttt{protected}, and \texttt{public}.

  \item \HMarks \textbf{[Program]} Demonstrate nested classes:
    create an outer class with one static nested class and one inner class; show object creation and member access.

  \item \HMarks \textbf{[Explain]} In 10--14 lines, explain why interfaces are needed in Java with a real use case.

  \item \HMarks \textbf{[Program]} Write a program that uses an anonymous class to implement a 1-method interface and executes it.

  \item \HMarks \textbf{[Program]} Write code showing that \texttt{finally} runs even when \texttt{return} is executed in \texttt{try}.
    Include output.
\end{enumerate}

\vfill
\noindent\textit{End of Assessment}

\end{document}

